\begin{exercise}{142}
Let $p\colon K\to\mathcal C$ be a diagram and let $q\colon I\to\mathcal C_{p/}$ be a further diagram. Let
$\bar q\colon K\star I\to\mathcal C$ be the associated map.\footnote{The book writes
$\bar q\colon I\star K\to\mathcal C$ here. The join--slice adjunction for
$q\colon I\to\mathcal C_{p/}$ instead gives a map $K\star I\to\mathcal C$.}
Show that there is an isomorphism $(\mathcal C_{p/})_{q/}\cong\mathcal C_{\bar q/}$.
Likewise, show that there is an isomorphism $(\mathcal C_{/p})_{q/}\cong(\mathcal C_{q'/})_{/p'}$ where $q'$ is the restriction
of $\bar q\colon I\star K\to\mathcal C$\footnote{The book writes
$q\colon I\star K\to\mathcal C$ here. In this second assertion,
$q\colon I\to\mathcal C_{/p}$ has transpose $\bar q\colon I\star K\to\mathcal C$;
since $q$ has domain $I$, the restriction is of $\bar q$.}
to $I$ and $p'$ is adjoint to $\bar q$.

\begin{proof}
Use the join--slice adjunctions of Corollary~1.4.17.
An $n$-simplex of $(\mathcal C_{p/})_{q/}$ is a map
$I\star\Delta^n\to\mathcal C_{p/}$ extending $q$, equivalently a map
$K\star(I\star\Delta^n)\to\mathcal C$ extending $\bar q$ on $K\star I$.
Join associativity identifies this with an $n$-simplex of
$\mathcal C_{\bar q/}$.

For the second assertion, take $q\colon I\to\mathcal C_{/p}$.
Its transpose is $\bar q\colon I\star K\to\mathcal C$;
write $q'=\bar q|_I$ and let $p'\colon K\to\mathcal C_{q'/}$
be its other transpose.
An $n$-simplex of $(\mathcal C_{/p})_{q/}$ is equivalently a map
$(I\star\Delta^n)\star K\to\mathcal C$ restricting to $\bar q$ on $I\star K$.
An $n$-simplex of $(\mathcal C_{q'/})_{/p'}$ is equivalently a map
$I\star(\Delta^n\star K)\to\mathcal C$ with the same restriction.
Associativity gives the required bijections.
In both assertions these bijections commute with every face and
degeneracy map, by naturality of the adjunctions and the join associator.
They therefore give the asserted canonical isomorphisms of simplicial sets.
\end{proof}
\end{exercise}
